\section{课程定位：为什么今天要讨论 ``可预测噪声''}

\subsection{从 ``分数崇拜'' 到 ``分数校准''}
Sida Wang 这讲的起点很直接：当前社区对 benchmark 的使用方式，经常默认 ``分数差异 = 能力差异''。在小规模任务上这个假设有时成立，但在大模型和 agentic workload 下，分数本身已经变成高噪声统计量。也就是说，榜单并非无价值，而是必须先回答 ``这个分数的误差有多大''、``这个差异是否统计显著''。

\begin{importantbox}{本讲核心命题}
Benchmark score 不是确定值，而是随机变量。只有把它放回概率与统计框架，排行榜才有可解释性。
\end{importantbox}

\subsection{噪声为何在 LLM 时代被放大}
传统监督学习评测里，输入固定、输出空间较小、评分规则清晰，因此噪声通常主要来自采样误差。LLM 评测不同：模型输出自由度更高，prompt 模板选择空间更大，解码策略和 judge pipeline 都会注入额外方差。尤其在 open-ended 任务和 agent benchmark 中，轨迹级交互使 variance 成倍上升。

\begin{knowledgebox}{LLM benchmark 噪声放大的三个结构性原因}
\begin{itemize}
    \item \textbf{生成随机性}：temperature、top-p、sampling seed 会改变输出分布。
    \item \textbf{协议随机性}：prompt wording、system instruction、few-shot selection 会改写任务难度。
    \item \textbf{评审随机性}：LLM-as-judge、人评一致性、rubric 模糊性都会造成测量噪声。
\end{itemize}
\end{knowledgebox}

\subsection{本章小结}
这讲不是反对 benchmark，而是反对把 noisy metric 当作 deterministic truth。后续所有技术点都围绕同一个目标：把 ``分数'' 变成 ``可校准、可比较、可复现实验信号''。

